\textbf{结论名称：}直线与圆的位置关系判定（d与r比较）

\textbf{适用条件：}已知圆的标准方程（或圆心、半径）及直线方程，且直线方程化为一般式 \(Ax+By+C=0\)，圆心到直线的距离 \(d\) 可求。

\textbf{变量说明：}设圆 \(C\) 的圆心为 \(P(x_0,y_0)\)，半径为 \(r\;(r>0)\)；直线 \(l:Ax+By+C=0\)；圆心到直线的距离为
\[
d = \frac{|Ax_0+By_0+C|}{\sqrt{A^2+B^2}}.
\]

\textbf{核心公式：}
\[
\boxed{
\begin{array}{c}
d<r \;\Longleftrightarrow\; \text{直线与圆相交（两个交点）},\\
d=r \;\Longleftrightarrow\; \text{直线与圆相切（一个交点）},\\
d>r \;\Longleftrightarrow\; \text{直线与圆相离（无交点）}.
\end{array}
}
\]

\textbf{使用场景：}快速判断直线与圆的位置关系；求切线方程或切点坐标；由位置关系反推参数取值范围；求弦长（相交时弦长 \(=2\sqrt{r^2-d^2}\)）。

\textbf{简单例子：}已知圆 \(x^2+y^2=4\)，直线 \(l:x-y+2\sqrt{2}=0\)。圆心 \((0,0)\)，半径 \(r=2\)。计算距离
\[
d=\frac{|0-0+2\sqrt{2}|}{\sqrt{1^2+(-1)^2}}=\frac{2\sqrt{2}}{\sqrt{2}}=2=r,
\]
因此直线与圆相切。